%% kapa: exam=1 round=35 year=2024 period=2 form=A subject=appraisal_law subject_ko=감정평가관계법규 seq=40 number=40
%% source: Q-Net 감정평가사 기출문제 (cst00309) artlSeq=5213229 file=raw/1008/5213229/2교시 문제지 원본.pdf
%% conversion: kordoc 4.24.51 -> tools/gen1.py (auto); review=proofread-sonnet; flags=none
\begin{question}{40}
도시 및 주거환경정비법령상 재건축사업의 관리처분계획에 관한 설명으로 옳은 것은?
\begin{choices}
  \item 관리처분계획에 포함될 분양대상자별 분양예정인 건축물의 추산액을 평가하기 위하여 시장ㆍ군수등이 선정ㆍ계약한 감정평가법인등을 변경하는 경우에는 조합총회의 의결을 거치지 않아도 된다.
  \item 토지등소유자에 대한 사업시행자의 매도청구에 대한 판결에 따라 관리처분계획을 변경하는 경우에는 시장ㆍ군수등의 변경인가를 받아야 한다.
  \item 사업시행자는 관리처분계획이 인가ㆍ고시된 날부터 90일 이내에 분양신청을 하지 않은 자와 손실보상에 관한 협의를 하여야 한다.
  \item 관리처분계획에 포함되는 세입자별 손실보상을 위한 권리명세 및 그 평가액은 시장ㆍ군수등이 선정한 2인 이상의 감정평가법인등이 평가한 금액을 산술평균하여 산정한다.
  \item 시장ㆍ군수등이 직접 관리처분계획을 수립하는 경우에는 토지등소유자의 공람 및 의견청취절차를 생략할 수 있다.
\end{choices}
\end{question}
\answer{1}
